\rtmtight
\begin{tblr}{width=\textwidth,colspec={|Q[c,m,2.1cm]|X[l,m]|},hlines,vlines,hspan=minimal,rowsep=1.5pt,colsep=3pt,row{1}={bg=headgray,font=\bfseries}}
\SetCell{c,m}\hfil\logo\hfil & \textbf{POLITEKNIK NEGERI MALANG}\par\textbf{JURUSAN TEKNOLOGI INFORMASI}\par\textbf{PROGRAM STUDI: D4 TEKNIK INFORMATIKA} \\
\SetCell[c=2]{l} \textbf{RENCANA TUGAS MAHASISWA (RTM) DAN RUBRIK PENILAIAN} & \\
Nama Mata Kuliah & Business Intelligence \\
Kode Mata Kuliah & RTI255005 \quad \textbf{sks / Jam perminggu:} 2 SKS/4 jam \quad \textbf{Semester:} 5 \\
Dosen Pengampu & Tim Pengajar Program Studi D4 TI \\
\end{tblr}
\assessmentblock{Tugas Data Warehouse dan ETL}{Case Method dan analisis data}{SCPMK0705-04101, SCPMK0706-04102, SCPMK1004-04103}{Mahasiswa menyelesaikan tugas merancang data warehouse dan ETL, mengimplementasikan OLAP dengan tools BI, serta memvisualisasikan insight bisnis yang relevan.}{Menganalisis kebutuhan data, merancang skema DW, mengimplementasikan proses ETL, membangun visualisasi, dan mendokumentasikan hasil.}{Skema data warehouse, skrip ETL, dashboard visualisasi, laporan, dan/atau bahan presentasi.}{Rancangan data warehouse dan skema ETL & 12 & Skema bintang/snowflake dirancang lengkap, proses ETL mencakup extract-transform-load dengan validasi, dan hasilnya konsisten dengan kebutuhan analitik bisnis. & Skema DW dan ETL sebagian besar benar, tetapi beberapa proses transformasi atau relasi dimensi-fakta belum tepat. & Skema DW tidak mencerminkan prinsip data warehouse atau proses ETL tidak dapat dijalankan. \\ \\
Implementasi OLAP dan tools BI & 10 & OLAP diimplementasikan dengan operasi slice/dice/drill-down/roll-up yang benar menggunakan tools BI (Power BI/Tableau/Metabase) secara efektif. & Implementasi OLAP berjalan tetapi beberapa operasi atau pemilihan tools belum optimal. & OLAP tidak diimplementasikan atau tools BI tidak digunakan secara fungsional. \\ \\
Kualitas visualisasi dan insight bisnis & 8 & Visualisasi memilih jenis chart yang tepat, menampilkan KPI yang relevan, dan insight yang disajikan dapat mendukung keputusan bisnis secara langsung. & Visualisasi cukup baik tetapi pemilihan chart atau relevansi insight belum sepenuhnya mendukung keputusan bisnis. & Visualisasi tidak tepat, tidak relevan, atau tidak menghasilkan insight bisnis yang bermakna. \\}

\assessmentblock{Kuis 1 Data Warehouse dan ETL}{Kuis}{SCPMK0705-04101}{Kuis tertulis untuk mengukur pemahaman konsep data warehouse, skema bintang/snowflake, dan proses ETL.}{Menjawab soal pilihan ganda dan/atau uraian terkait konsep DW, ETL, dan OLAP.}{Jawaban tertulis kuis.}{Pemahaman konsep data warehouse & 3 & Konsep DW, skema bintang, dimensi-fakta, dan peran ETL dijelaskan dengan benar dan contoh penerapannya tepat. & Sebagian besar konsep dipahami tetapi contoh atau penjelasan masih kurang lengkap. & Konsep DW atau ETL tidak dipahami atau contoh penerapan keliru. \\ \\
Analisis skema dan proses ETL & 2 & Kasus ETL dianalisis dengan tepat: tahap extract, transform, dan load diidentifikasi benar sesuai skenario bisnis. & Sebagian besar analisis benar tetapi beberapa tahap ETL atau relasi skema masih keliru. & Analisis tidak tepat atau tidak menunjukkan pemahaman proses ETL. \\}

\assessmentblock{UTS Data Warehouse dan OLAP}{UTS tertulis dan analisis}{SCPMK0705-04101, SCPMK0706-04102}{Evaluasi tengah semester untuk mengukur kemampuan merancang data warehouse, proses ETL, dan mengimplementasikan analisis OLAP menggunakan tools BI.}{Mengerjakan soal analisis kasus, merancang skema DW, dan menjelaskan implementasi OLAP.}{Dokumen jawaban, rancangan skema, dan analisis OLAP.}{Implementasi dan analisis data warehouse & 7 & Data warehouse dirancang dengan skema yang tepat, proses ETL dapat dijelaskan dengan benar, dan analisis kebutuhan data bisnis dirumuskan dengan akurat. & Rancangan DW sebagian besar benar tetapi analisis kebutuhan atau detail ETL masih belum lengkap. & Rancangan DW tidak mencerminkan prinsip data warehouse atau analisis tidak dapat dilakukan. \\ \\
Perancangan dan analisis OLAP & 6 & OLAP diimplementasikan dengan operasi multidimensional yang benar dan hasilnya dapat diinterpretasikan untuk mendukung keputusan bisnis. & OLAP sebagian besar benar tetapi interpretasi hasil atau pemilihan operasi masih belum optimal. & OLAP tidak diimplementasikan atau hasilnya tidak dapat diinterpretasikan. \\}

\assessmentblock{Kuis 2 Visualisasi dan Dashboard BI}{Kuis}{SCPMK1004-04103}{Kuis tertulis untuk mengukur pemahaman prinsip visualisasi data dan evaluasi kualitas dashboard BI.}{Menjawab soal terkait prinsip visualisasi, pemilihan jenis chart, dan evaluasi kualitas dashboard.}{Jawaban tertulis kuis.}{Pemahaman prinsip visualisasi data & 3 & Prinsip visualisasi (clarity, accuracy, efficiency) dan pemilihan jenis chart yang sesuai konteks bisnis dijelaskan dengan benar. & Sebagian besar prinsip dipahami tetapi penerapan pada kasus spesifik masih kurang tepat. & Prinsip visualisasi tidak dipahami atau pemilihan chart tidak relevan. \\ \\
Kemampuan evaluasi kualitas dashboard & 2 & Kualitas dashboard dievaluasi berdasarkan keterbacaan, relevansi KPI, dan kemampuan mendukung keputusan bisnis secara tepat. & Evaluasi dashboard cukup baik tetapi beberapa dimensi kualitas belum dipertimbangkan. & Evaluasi tidak tepat atau tidak menggunakan kriteria kualitas yang relevan. \\}

\assessmentblock{PBL Dashboard BI}{Project Based Learning}{SCPMK0706-04102, SCPMK1004-04103}{Mahasiswa membangun dashboard BI interaktif yang mengintegrasikan analisis OLAP, data mining, dan visualisasi data untuk menghasilkan insight bisnis yang bermakna.}{Menganalisis kebutuhan bisnis, merancang arsitektur BI, mengimplementasikan dashboard, menguji fungsionalitas, dan mempresentasikan insight.}{Dashboard BI fungsional, laporan analisis, dokumentasi teknis, dan bahan presentasi.}{Integrasi OLAP dan analisis data & 8 & Dashboard mengintegrasikan OLAP dengan operasi multidimensional yang benar, data mining digunakan untuk menghasilkan insight prediktif, dan alur data dari DW ke dashboard konsisten. & Integrasi OLAP dan data mining sebagian besar benar tetapi konsistensi alur data atau analisis prediktif masih terbatas. & OLAP tidak terintegrasi atau dashboard tidak dapat menghasilkan analisis data yang bermakna. \\ \\
Kualitas dashboard dan visualisasi & 9 & Dashboard interaktif dengan visualisasi yang tepat, KPI yang relevan, storytelling data yang efektif, dan performa yang baik. & Dashboard berfungsi tetapi interaktivitas, pemilihan visualisasi, atau storyline data masih kurang optimal. & Dashboard tidak interaktif, visualisasi tidak tepat, atau tidak dapat menyampaikan insight bisnis. \\ \\
Dokumentasi, insight, dan presentasi & 5 & Insight bisnis dirumuskan secara akurat dari data, terdokumentasi dengan baik, dan dipresentasikan dengan meyakinkan kepada pemangku kepentingan. & Insight cukup relevan tetapi dokumentasi atau penyajian kepada pemangku kepentingan masih belum kuat. & Insight tidak akurat, tidak terdokumentasi, atau tidak dapat dipresentasikan secara meyakinkan. \\}

\assessmentblock{UAS Presentasi dan Demo Dashboard BI}{UAS presentasi dan demo}{SCPMK0706-04102, SCPMK1004-04103}{Mahasiswa mempresentasikan dan mendemonstrasikan dashboard BI yang telah dibangun, menjelaskan insight bisnis yang dihasilkan, dan mempertahankan keputusan desain analitik.}{Melakukan demo dashboard, mempresentasikan insight bisnis, menjawab pertanyaan, dan menjelaskan keputusan teknis dan analitik.}{Dashboard BI final, laporan akhir, dan bahan presentasi.}{Kelengkapan dan fungsionalitas dashboard & 10 & Dashboard berjalan lengkap dengan semua fitur yang diminta, OLAP terintegrasi, visualisasi responsif, dan koneksi data berfungsi dengan benar. & Sebagian besar fitur berfungsi tetapi beberapa integrasi OLAP atau visualisasi masih ada kekurangan. & Dashboard tidak berjalan atau banyak fitur yang tidak fungsional. \\ \\
Akurasi analisis dan insight bisnis & 9 & Insight yang disajikan akurat, berbasis data, relevan untuk pengambilan keputusan strategis, dan dapat dipertahankan dengan argumen analitik yang kuat. & Insight cukup relevan tetapi akurasi analisis atau argumentasi analitik masih perlu diperkuat. & Insight tidak akurat, tidak berbasis data, atau tidak dapat dipertahankan secara analitik. \\ \\
Kualitas presentasi dan defense & 6 & Presentasi terstruktur, demo berjalan lancar, pertanyaan dijawab dengan tepat, dan keputusan desain analitik dijelaskan dengan meyakinkan. & Presentasi cukup baik tetapi demo atau jawaban atas pertanyaan masih belum optimal. & Presentasi tidak terstruktur, demo gagal, atau tidak dapat menjawab pertanyaan terkait keputusan desain. \\}

\begin{tblr}{width=\textwidth,colspec={|Q[l,m,3.2cm]|X[l,m]|},hlines,vlines,hspan=minimal,row{1-3}={bg=headgray},rowsep=1.5pt,colsep=3pt}
Jadwal Pelaksanaan: & Minggu 1--17 sesuai RPS. Setiap evaluasi memiliki RTM dan rubrik multi-kriteria. \\
Lain-Lain Yang Diperlukan: & LMS, repositori/data, perangkat lunak BI sesuai mata kuliah, dan perangkat presentasi. \\
Pustaka: & Mengacu pustaka utama dan pendukung pada bagian RPS. \\
\end{tblr}
